% Terminal apparatus of the concise study: the scope note and the References
% for the sources this companion actually uses. The revision timestamp and the
% rights colophon follow from the shared generation record.

\section{Appendix: Scope and Qualifications}

{\small
\textbf{Edition, formulary and occurrence.} This concise study treats the Mass
of the Twenty-fifth Sunday in Ordinary Time, Year~A, as the \work{Roman
Missal}, Third Edition, and the \work{Lectionary for Mass}, no.~133, give it
for the dioceses of the United States, in the national calendar alone, for the
occurrence of 20~September 2026. The antiphons and orations serve Years A, B
and~C; the readings, psalm and Alleluia verse are Year~A's. Both
Communion antiphons are treated.

\textbf{Relation to the expansive study.} Every claim here is drawn from the
canonical study of this same formulary, which prints the appointed Scripture
and develops the three interpretations at length, each with its own literal,
allegorical, moral and anagogical senses; nothing evidence-dependent is added.
The four rows on the first page orient a reader and do not replace those
distillations. The second page reprints the study's dossier with the dates,
relation labels, alternatives and unresolved states exactly as the generated
chronology record gives them.

\textbf{Text and English.} The Missal formulary was read for the study in the
Latin of the \work{Missale Romanum}, editio typica tertia (2002), in a
digitally typeset reproduction; no United States altar book was inspected, so
the approved English orations were not read in any witness. The Lectionary's
boundaries rest on the national calendar for 2026, the bishops' conference's
dated page of readings and no.~133 of the Latin \work{Ordo lectionum Missae}
(1981); the printed United States Lectionary was not inspected. No critical apparatus of the Greek New
Testament was opened, and nothing is said here about which manuscripts carry
the second sentence of Matthew 20:16 or why the appointment is lettered 16a.
All Scripture is quoted in the Douay--Rheims version as revised by Bishop
Challoner, a public-domain study translation of the Clementine Vulgate. It is
not the text proclaimed in the celebration: the controlling English is that of
the \work{Lectionary for Mass for Use in the Dioceses of the United States of
America} and, for the antiphons and orations, that of the \work{Roman Missal},
Third Edition, neither of which is reproduced. The orations and the Entrance
Antiphon are identified by Latin incipit and described; the descriptions are
not translations and may not be quoted as the prayers' words. English glosses
of patristic Latin not drawn from a named translation render the Latin beside
them.

\textbf{Reception and its limits.} The witnesses named are those the study's
bounded search read at the loci in the References. Chrysostom was read in the
English of the Nicene and Post-Nicene Fathers, as were Augustine's \work{Sermo}
87, his tractate on John and his work on the predestination of the saints;
Gregory in an electronic Latin text; Jerome, Bede, Hilary and Aquinas on page
images of the printings named, except that of Aquinas on Matthew only p.~262
and two sentences of pp.~258--261 were read on the page, so that what is
reported from those four pages is paraphrased. Of Augustine's
\work{Enarrationes}, 118, s.~4.2 and 144.11 were read on the columns of Migne,
and 144.22 rests on the electronic text of augustinus.it alone. Bellarmine was
read in a transcription of O'Sullivan's abridged English of 1866, and the Latin
was not compared. No witness was collated with a modern critical edition.
Nothing is said here about who first proposed either reading of the hours or about the order in time or
the mutual dependence of the witnesses. No commentary on the orations or on
the Entrance Antiphon was sought, and no witness cited comments on this Mass.
The liturgical books state three relations among its elements
(\work{Praenotanda} 106; \work{General Instruction of the Roman Missal} 61,
62), and that the Alleluia verse's basis and the second reading both concern
Philippi is an observation about the texts. Every other connection between one
element and another is synthesis, grounded in a witness's statement about his
own passage where one is cited and otherwise resting on the wording of the
texts; none claims that the formulary was composed to make it.

\textbf{Chronology, rights and records.} Every date on the second page comes
from the Scripture chronology corpus through the record generated for this
formulary. For the two psalms that record holds only the modern critical
boundary of the Psalter, summarized from the introduction to the Psalms in the
\work{New American Bible, Revised Edition}, and no traditional date; for
Isaiah, Matthew, John, Acts and Philippians it holds the dates reported in the
\work{Catholic Encyclopedia}, and later critical positions are named without
figures. The Gospel's narrated event is undated. The Scripture, patristic and
encyclopedia texts quoted are in the public domain, Gregory's Latin being
quoted in short phrases from a licensed electronic edition; the liturgical
books, the \work{Nova Vulgata} and the New American Bible introduction are
protected and are cited, with only incipits, the three titles of the
\work{Ordo} and a few words of the Latin books quoted. The audits are \texttt{propers/verified.md}, \texttt{research/scope.md},
\texttt{research/interpretations.md} and, for this companion's production,
\texttt{research/production-review.md}.
\par}

\section*{References}
% A starred heading sets no mark, so without this the Appendix above would go
% on heading the pages of the bibliography.
\runninghead{References}

{\small
% One list in three runs: liturgical books and Bibles; Fathers, Doctors and
% saints; chronology sources.
\begin{itemize}\setlength{\itemsep}{0.15em}\setlength{\parskip}{0pt}
\item \work{Lectionary for Mass for Use in the Dioceses of the United States of America}, second typical edition, vol.~I, no.~133, as witnessed by the United States Conference of Catholic Bishops, \work{Liturgical Calendar} 2026, pp.~5 and 38, and its page of readings for 20~September 2026, \url{https://bible.usccb.org/bible/readings/092026.cfm}; \work{General Instruction of the Roman Missal}, United States edition of 2011 with the emendations of 2021, nos.~61--62.
\item \work{Missale Romanum}, editio typica tertia (Typis Vaticanis, 2002), \latin{Dominica XXV ``per annum''}, and the rubrics of \latin{Tempus per annum}, no.~6.
\item \work{Ordo lectionum Missae}, editio typica altera (Libreria Editrice Vaticana, 1981): \work{Praenotanda} 105--107; no.~133, \latin{Dominica vigesima quinta}, Year~A.
\item The Holy Bible, Douay--Rheims version, revised by Bishop Richard Challoner (Project Gutenberg text): Lev 18:5; Ps 118:4--5; Ps 144; Isa 1:1; 55; Mt 19:27--20:19; 22:40; Jn 10:11--16; Acts 16:12--15; Rom 10:5; Phil 1; Col 4:3.
\item \work{Biblia Sacra Vulgatae editionis} (Clementine Vulgate), eBible.org text: Lev 18:5; Ps 118:4--5; Mt 20:16; Jn 10:14; Rom 10:5. \work{Nova Vulgata Bibliorum Sacrorum editio}, as published by the Holy See at vatican.va: Mt 20:16.
\item St John Chrysostom, \work{Homilies on Matthew} 64, trans. G. Prevost, rev. M.~B. Riddle, Nicene and Post-Nicene Fathers, 1st ser., vol.~10 (the series in the CCEL electronic text); \work{Homilies on Philippians} 3--4, 1st ser., vol.~13; \work{Homilies on Acts} 35, 1st ser., vol.~11.
\item St Augustine, \work{Sermo} 87 (Sermon 37 on the New Testament Lessons), Nicene and Post-Nicene Fathers, 1st ser., vol.~6, sections 1, 4--11; \work{De praedestinatione sanctorum} 20.41, 1st ser., vol.~5; \work{In Iohannis evangelium tractatus} 47.5, 1st ser., vol.~7.
\item St Augustine, \work{Enarrationes in Psalmos} 118, s.~4.2, and 144.11 (Migne, PL 37, Paris, 1845, cols.~1510 and 1876); 144.22 in the Latin text of \url{https://www.augustinus.it/latino/esposizioni_salmi/}.
\item St Gregory the Great, \work{Homiliae in Evangelia} 14.1--5 and 19.1--4, Latin (Wikisource text).
\item St Jerome, \work{Commentarii in Matthaeum} III, on 20:1--16 (Migne, PL 26, Paris, 1845, cols.~140--142); \work{Commentarii in Isaiam} XV, on 55:6--11 (Migne, PL 24, Paris, 1865, cols.~553--555).
\item St Hilary of Poitiers, \work{Tractatus super Psalmos}, ed. A. Zingerle, CSEL 22 (Vienna, 1891): on Ps~118, Aleph 12 (p.~367); on Ps~144, 14 (p.~834).
\item St Bede, \work{Liber retractationis in Actus apostolorum}, c.~16 (Migne, PL 92, Paris, 1862, col.~1025).
\item St Thomas Aquinas, \work{Super Evangelium Matthaei} c.~20, lect.~1, and \work{Super Evangelium Ioannis} c.~10, lect.~4, in \work{Opera}, editio altera Veneta, vol.~3 (Venice: Bettinelli, 1745), pp.~258--262 and 662; \work{Super epistolam ad Philippenses} c.~1, lect.~3, in \work{In omnes S.~Pauli Apostoli epistolas commentaria}, vol.~2 (Li\`ege: Dessain, 1857), pp.~371--372; \work{Expositio super Isaiam ad litteram} c.~55, in \work{Opera omnia}, vol.~14 (Parma: Fiaccadori, 1863), p.~556.
\item St Robert Bellarmine, \work{A Commentary on the Book of Psalms}, trans. J. O'Sullivan (1866, abridged), in the eCatholic2000 transcription: on Ps~144:17.
\item \work{The Catholic Encyclopedia} (New York, 1907--1912), New Advent text: A.~E. Breen, ``Acts of the Apostles'' (1907); J. Howlett, ``Biblical Chronology'' (1908); C. Souvay, ``Isaias'' (1910); L. Fonck, ``Gospel of St.~John'' (1910); J.~E. Jacquier, ``Gospel of St.~Matthew'' (1911); F. Prat, ``St.~Paul'' (1911); A. Vander Heeren, ``Epistle to the Philippians'' (1911); A. Durand, ``The New Testament'' (1912). United States Conference of Catholic Bishops, \work{New American Bible, Revised Edition}, introduction to the Psalms (summarized, not quoted).
\end{itemize}
\par}
